\bookchapter{Capitolo 53: Conservare i dati e parlare con la rete}

\chapterintro{Un'app che serve solo a «guardare» è morta; un'app che salva dati e si collega alla rete è viva. In questo capitolo conosci il «trio» del livello dati di Android: DataStore (le piccole configurazioni), Room (i dati strutturati), Retrofit (chi parla con il server), più il loro «direttore d'orchestra» — le coroutine. A capitolo finito potrai scrivere app vere, capaci di ricordare le abitudini dell'utente e di aggiornare i contenuti in tempo reale.}

\booksection{Come scegliere dove salvare i dati: prima guarda questa tabella}

\begin{center}
\begin{tabularx}{\linewidth}{|X|X|X|X|}
\hline
\textbf{Soluzione} & \textbf{Adatta a salvare cosa} & \textbf{Caratteristiche} & \textbf{Dove si usa in questo libro} \tabularnewline
\hline
DataStore & Configurazioni sparse (tema, interruttori) & Coppie chiave-valore, semplice e sicuro & Le impostazioni del registro delle spese, capitolo 54 \tabularnewline
Room & Dati strutturati (voci di spesa, utenti) & Una vera incapsulazione del database SQLite, interrogabile & Le voci del registro, capitolo 54 \tabularnewline
Retrofit & Dati del server (non si salvano, si scaricano dalla rete) & Lo standard di facto per le richieste HTTP & I dati meteo, capitolo 55 \tabularnewline
\hline
\end{tabularx}
\end{center}

\begin{center}\small\color{gray}La ripartizione dei compiti del trio dei dati\end{center}

Impara questa filastrocca: «le piccole configurazioni in DataStore, i dati in blocco in Room, i dati fuori porta passano da Retrofit». Le SharedPreferences, comuni nei vecchi tutorial, sono state sostituite da DataStore, e il SQLite scritto a mano da Room — impara direttamente il nuovo, non riprendere la vecchia strada.

\booksection{DataStore: salvare le piccole configurazioni}

Per ricordare valori sparsi come «la valuta scelta l'ultima volta dall'utente» o «se la modalità scura è attiva», usa DataStore (versione Preferences):

\begin{lstlisting}[language=Kotlin]
val Context.dataStore by preferencesDataStore(name = "settings")

// Scrittura: salva un valore
suspend fun saveTheme(isDark: Boolean) {
    context.dataStore.edit { prefs ->
        prefs[THEME_KEY] = isDark
    }
}

// Lettura: preleva (il Flow notifica i cambiamenti da solo)
val themeFlow: Flow<Boolean> = context.dataStore.data
    .map { it[THEME_KEY] ?: false }
\end{lstlisting}

Nota che viene restituito un Flow — un flusso di dati che «scorre». Appena il valore cambia, tutti i punti che lo osservano ricevono da soli il nuovo valore. Questo concetto (la reattività) è la pietra angolare dell'Android moderno: lo metterai in pratica nel capitolo 54.

\booksection{Room: dare ai dati una «casa vera»}

Un registro delle spese ha centinaia o migliaia di voci, e deve interrogare e riepilogare per categoria e per mese — serve un database. Room è l'incapsulamento ufficiale di SQLite e rende «creare tabelle e interrogare i dati» naturale quanto scrivere Kotlin. Servono tre cose:

\begin{lstlisting}[language=Kotlin]
// 1. Entity: la struttura della tabella (una tabella = un data class)
@Entity(tableName = "records")
data class Record(
    @PrimaryKey(autoGenerate = true) val id: Long = 0,
    val amount: Double,
    val category: String,
    val note: String,
    val date: String
)

// 2. Dao: l'interfaccia delle operazioni (CRUD)
@Dao
interface RecordDao {
    @Query("SELECT * FROM records ORDER BY date DESC")
    fun getAll(): Flow<List<Record>>

    @Insert
    suspend fun insert(record: Record)

    @Delete
    suspend fun delete(record: Record)

    @Query("SELECT SUM(amount) FROM records WHERE category = :c")
    suspend fun sumByCategory(c: String): Double
}

// 3. Database: assemblaggio finale
@Database(entities = [Record::class], version = 1)
abstract class AppDatabase : RoomDatabase() {
    abstract fun recordDao(): RecordDao
}
\end{lstlisting}

Questo modello a «tre pezzi» si può copiare identico in qualunque tuo progetto — basta cambiare i nomi dei campi. Le query restituiscono Flow e l'interfaccia si aggiorna da sola al passo con i dati: una meraviglia.

\booksection{Retrofit: parlare con il server}

Meteo, notizie, tassi di cambio... questi dati non stanno nel telefono, stanno sui server degli altri. Retrofit rende «chiedere dati al server» semplice come chiamare una funzione. Tre passi:

\begin{lstlisting}[language=Kotlin]
// 1. Definisci l'interfaccia: descrivi le richieste HTTP come funzioni
interface WeatherApi {
    @GET("v1/weather")
    suspend fun getWeather(@Query("city") city: String): WeatherResponse
}

// 2. Crea l'istanza
val api = Retrofit.Builder()
    .baseUrl("https://api.example.com/")
    .addConverterFactory(GsonConverterFactory.create())
    .build()
    .create(WeatherApi::class.java)

// 3. Chiamala come una funzione (in una coroutine)
val result = api.getWeather("Pechino")
\end{lstlisting}

\textbf{Attenzione alle trappole: le due regole d'oro della rete}\par

(1) Le richieste di rete non si fanno mai nel thread principale — congela l'interfaccia e può far crashare l'app (NetworkOnMainThreadException). Tutte le chiamate di rete vanno messe in una coroutine (viewModelScope.launch). (2) L'HTTP in chiaro è vietato per impostazione predefinita: da Android 9 in poi serve l'HTTPS. Se per i test devi usare un indirizzo locale in http://, devi dichiarare esplicitamente usesCleartextTraffic nel Manifest — ma un prodotto vero deve viaggiare tutto in HTTPS.

\booksection{Il parsing del JSON: «tradurre» i dati di rete}

I dati che restituisce il server sono per lo più JSON — una fila di testo che sembra un «dizionario annidato». Retrofit ha bisogno di un «interprete» che li trasformi in oggetti Kotlin. Le due soluzioni più diffuse nel 2026:

\begin{center}
\begin{tabularx}{\linewidth}{|X|X|X|}
\hline
\textbf{Soluzione} & \textbf{Caratteristiche} & \textbf{Quando consigliarla} \tabularnewline
\hline
Gson / Moshi & La soluzione classica basata sulla riflessione, a costo di configurazione zero & Piccoli progetti partiti in fretta, con pochi campi \tabularnewline
kotlinx.serialization & La soluzione ufficiale a compile-time, più veloce e più sicura & Prima scelta per i nuovi progetti (la corrente principale del 2026) \tabularnewline
\hline
\end{tabularx}
\end{center}

\begin{center}\small\color{gray}Le due grandi soluzioni per il parsing del JSON\end{center}

\begin{lstlisting}[language=Kotlin]
// Aggiungi l'annotazione @Serializable al data class
@Serializable
data class WeatherResponse(
    val weather: List<Weather>,
    val main: Main,
    val name: String
)

// Retrofit riceve il Converter corrispondente (supporto ufficiale da Retrofit 2.11+):
// .addConverterFactory(json.asConverterFactory("application/json".toMediaType()))
\end{lstlisting}

Con qualunque soluzione la regola d'oro è una: prima guarda il JSON realmente restituito, poi definisci i campi. Tanti bug di «parsing fallito» nascono da nomi di campo sbagliati o livelli di annidamento indovinati a caso — stampa il JSON originale e confrontalo: in dieci secondi trovi il punto esatto.

\booksection{Le coroutine: la «lingua franca» della concorrenza in Android}

\booksection{Le coroutine: la «lingua franca» della concorrenza in Android}

I suspend e launch che compaiono più volte sopra vengono dalle coroutine di Kotlin (Coroutine) — la soluzione di concorrenza ufficiale di Android. Puoi capirlo alla buona così: far sì che le operazioni che «chiedono di aspettare» (rete, database) non blocchino l'interfaccia.

\begin{lstlisting}[language=Kotlin]
// Avvia una coroutine nel ViewModel
viewModelScope.launch {
    val records = dao.getAll().first()  // suspend: si sospende e attende da solo
    // aggiorna lo stato dell'interfaccia...
}

// Promemoria delle parole chiave:
// suspend   — "devo aspettare, non bloccate gli altri"
// launch    — "Al lavoro!"
// viewModelScope — l'"officina" legata alla schermata; si chiude da sola quando la schermata viene distrutta
\end{lstlisting}

Se parti da zero, per ora basta ricordarne l'uso; il principio del «sospendere e riprendere» che c'è dietro potrai approfondirlo più avanti. Nel lavoro di tutti i giorni il 90\% di quello che scriverai è questo schema identico.

\begin{tcbpractice}{Mani all'opera: crea una mini app «una frase al giorno»}
Trova un'API gratuita di citazioni (per esempio hitokoto o simili), usa Retrofit per scaricare la «frase del giorno», mostrala al centro dello schermo e aggiungi un pulsante «Cambia frase» (che richiede di nuovo dentro una coroutine). Poi salva le frasi che ti piacciono in una «raccolta» con Room. Questo esercizio usa tutti assieme DataStore/Room/Retrofit/coroutine — finito quello, i capitoli 7 e 8 saranno un gioco da ragazzi.

\end{tcbpractice}

\begin{tcbsummary}{Riepilogo del capitolo}
\begin{itemize}
  \item Filastrocca: piccole configurazioni in DataStore, dati in blocco in Room, dati fuori porta via Retrofit
  \item DataStore restituisce un Flow: se il valore cambia, gli osservatori lo ricevono da soli
  \item Il trio di Room: @Entity (tabella) + @Dao (operazioni) + @Database (assemblaggio), modello riutilizzabile così com'è
  \item I tre passi di Retrofit: definisci l'interfaccia crea l'istanza la chiami in una coroutine
  \item Le due regole d'oro: la rete non va nel thread principale; un prodotto vero viaggia solo in HTTPS
  \item Coroutine = un'attesa che non blocca l'interfaccia; viewModelScope è l'officina riservata alla schermata
\end{itemize}

\end{tcbsummary}
